\begin{table*}[htbp]
\centering
\small
\caption{Future directions and recommendations derived from the extraction framework.}
\label{tab:recommendations}
\begin{tabularx}{\textwidth}{@{}>{\raggedright\arraybackslash}p{0.20\textwidth}YY@{}}
\toprule
\textbf{Recommendation} & \textbf{Priority} & \textbf{Evidence basis and suggested action} \\
\midrule
Standardize terminology & High & Report explicit criteria for digital model, digital shadow, digital twin, patient digital twin, virtual patient, and in silico patient. \\
Define maturity levels & High & Classify patient specificity, hierarchy, temporal updating, data coupling, feedback, validation, and clinical integration. \\
Strengthen validation & High & Require clear validation datasets, internal/external validation where relevant, and sensitivity or uncertainty analysis. \\
Build interoperable infrastructure & High & Report applicable FHIR, HL7, DICOM, OMOP, ontology, identifier, and governance mechanisms. \\
Provide reproducible supplements & High & Publish row-level extraction outputs, executable rules, data-quality checks, and plotting datasets. \\
Study workflow and governance & Medium & Evaluate clinician trust, liability, consent, bias, interpretability, equity, and implementation burden. \\
Move toward prospective evidence & High & Distinguish conceptual promise from retrospective validation, prospective studies, and deployed clinical workflows. \\
\bottomrule
\end{tabularx}
\end{table*}
